%% sections/05-discussion.tex

\section{Discussion}
\label{sec:discussion}

\subsection{The Amplification Mechanism}

The amplification mechanism we propose is: resource exhaustion, when legible and
emotionally architectured, \textit{forces character decisions to be visible}. In a
no-pressure session, a PC can resolve a situation through tactical default --- attack
until the encounter ends, spend a healing word when HP dips, use a spell slot when
the encounter demands it. Under resource pressure, tactical defaults become costly:
every spell slot spent is permanently gone; every HP taken persists into the next
session; every unit member lost will not return. The party must decide whether each
expenditure is worth its narrative cost.

This visibility mechanism works differently from the familiar difficulty-ratchet of
resource pressure in tactical games. In tactical games, resource pressure makes
decisions harder (fewer options) and forces triage (which problem to solve with limited
resources). In narrative games with emotionally architectured resources, pressure does
something additionally: it makes the party's values visible as resource decisions.
Conserving healing supplies is not just resource management; it is a decision about who
the party is willing to let go without help. Preserving unit members is not just
tactical force-maintenance; it is a decision about who the party considers worth
protecting at tactical cost.

Campaign 4's healer PC, Sava, illustrates this. Over 7 sessions, her spell-slot decisions
were explicitly shaped by her grief-paragraph pattern (not spending prematurely,
because early spending had cost someone in her backstory). The resource mechanic
made this grief-paragraph-driven behavior visible as a character decision across the
full campaign, rather than letting it dissolve into session-by-session optimization.
Without the between-session HP carry-over, Sava's slot decisions in Session 3 would
have no relationship to her slot availability in Session 7. With it, her Session 3
restraint was a Session 7 investment --- and the rubric could score it as such.

\subsection{Why the Cohesion Die Amplifies More Strongly}

The Cohesion Die's superior amplification in sessions 5--7 (character-primary vs.\
character-inflected tactical in panel assessment) is explained by three properties
distinguishing people-tracking from materials-tracking:

\begin{enumerate}
  \item \textbf{Named loss:} Each Cohesion Die degradation represents named unit members.
    The party knows who is gone --- not ``3 units lost'' but ``Carra, Fennet, and Voss
    are gone.'' Material depletion (healing supplies at d4) does not have names. Named
    loss activates what the rubric calls the \textit{grief-paragraph channel}: the party
    responds to loss of named individuals through character identity (grief, duty,
    commitment) rather than through tactical accounting.

  \item \textbf{Irreversibility:} Unit members who are gone do not return. Healing supplies
    at d4 could theoretically be resupplied through a side-quest or a lucky discovery;
    the 4 unit members lost between sessions 3 and 7 cannot be. Irreversibility
    increases the weight of each Cohesion Die decision in two ways: it raises the stakes
    of each pre-combat discussion (should we even engage?) and it makes post-combat
    accounting a character moment (who did we lose?).

  \item \textbf{Active preservation:} The party in Campaign 6 consistently made decisions
    that prioritized keeping unit members alive over tactical advantage. In Sessions 4--6,
    at least one PC per session made a tactically sub-optimal choice (measured by expected
    damage output or action economy) to preserve a unit member. This preservation behavior
    \textit{is itself} character-driven decision-making: it is the party choosing to
    sacrifice efficiency to protect specific named people. The Supply Die does not produce
    the same behavior; the party treated the Supply Die as a constraint to be managed, not
    as a responsibility to be honored.
\end{enumerate}

The party's relationship to the two resources is therefore qualitatively different.
Campaign 4 managed the Supply Die. Campaign 6 \textit{cared about} the Cohesion Die.
The distinction between managing a constraint and caring about what the constraint
represents is the design variable that produces the stronger amplification effect.

\subsection{Implications for Resource Mechanic Design}

For designers building resource mechanics for long-form narrative campaigns, the
findings suggest four design principles:

\begin{enumerate}
  \item \textbf{Make the resource legible.} A degrading die with named categories
    (d12 $\to$ d10 $\to$ d8 $\to$ d6 $\to$ d4 $\to$ exhausted) is more effective than
    HP numbers alone. The party must be able to see their cumulative depletion at each
    session without calculation. Legibility is the prerequisite for the resource becoming
    emotionally real.

  \item \textbf{Make what the resource represents matter.} Resources that represent
    people (or named relationships, or specific commitments) produce stronger amplification
    than resources that represent abstract materials. Name what the resource is tracking,
    not just what it controls. If the Cohesion Die is at d8, the party should know that
    it is at d8 because Carra and Voss and Fennet are gone --- not just because ``the
    unit has been degraded.''

  \item \textbf{Allow preservation as character expression.} Design the resource so
    that actively preserving it is a valid, if tactically sub-optimal, choice. If
    preservation is never rational --- if spending the resource is always the right
    move --- the party will spend it. The Cohesion Die's order-defiance degradation
    path (the die degrades if PCs defy orders without unit consensus) makes preservation
    a multi-layered commitment: protect members, but also maintain unit cohesion through
    legitimate decision-making. This second degradation path rewards character behavior
    (consulting unit members before committing) rather than just tactical behavior
    (keeping the die high).

  \item \textbf{Do not resupply prematurely.} The between-session no-long-rest mechanic
    is what made depletion cumulative and visible in both resource campaigns. If the
    party had been resupplied mid-campaign, the die would have reset, the grief-paragraph
    channel would have been interrupted, and the monotone-rising CA trajectory would not
    have accumulated. In a 7-session campaign, the resource should degrade for at least
    5 sessions before any recovery mechanism is permitted --- if recovery is permitted
    at all.
\end{enumerate}

\subsection{The Attrition Paradox}

The central finding --- that resource exhaustion amplifies rather than suppresses
character-driven decision-making --- resolves what we call the \textit{attrition paradox}
in long-form narrative game design. Conventional wisdom holds that resource pressure
crowds out narrative engagement: when survival is at stake, players optimize rather
than characterize. The Marathon data suggests this is true for short-form games
(where resource pressure arrives too suddenly to be architectured) but inverts in
long-form campaigns where the resource has had time to become emotionally freighted.

By Session 5 of Campaign 6, the Cohesion Die at d10 is not just a tactical number;
it is the accumulated record of who the party has lost and kept. Every decision that
touches the die is simultaneously a tactical decision and a character statement. The
party does not stop characterizing under pressure; they characterize \textit{through}
pressure, using resource decisions as the primary medium for expressing character
values. The attrition paradox dissolves once we recognize that the resource is not
competing with narrative --- it \textit{is} the narrative, made mechanical.

\subsection{Archetype Confound and Causal Inference}
\label{subsec:archetype-confound}

The Unit Cohesion Die and the military party archetype (duty-bound PCs, chain of command,
named unit members as social commitments) were introduced simultaneously in Campaign 6.
The $+3.0$ CA-point-session advantage observed in sessions 5--7 could therefore be
partly attributable to the duty-bound party archetype rather than to the resource mechanic
itself. A party whose PCs are designed around explicit duty and chain-of-command commitments
may produce character-primary decisions independently of any mechanical resource tracking.
This is the archetype confound in its sharpest form.

\paragraph{Available evidence against pure archetype confound.}
Three campaigns provide relevant partial controls:

\begin{itemize}
  \item \textbf{C5 (professional-code archetype, no resource mechanic):} Campaign 5
    runs a professional-code party (guild duty, contractual loyalty, craft identity) ---
    a duty-adjacent archetype with strong PC-level commitments but no degrading-resource
    mechanic. C5's CA trajectory is stable: 9.0 at Session 1, 9.1 at Session 7, with
    no monotone increase (see Table~\ref{tab:ca-trajectories}). A duty-adjacent party
    without a resource mechanic does not show the monotone-rising amplification pattern.
    This is the most direct available comparison for isolating archetype from mechanic.

  \item \textbf{C7 (stranger archetype, no resource mechanic):} Campaign 7 shows a
    late-rising trajectory, with the CA increase concentrated in sessions 4--7 as
    inter-PC relationships form. Relationship formation, not session 1 mechanics, drives
    the late-campaign rise in C7. Campaign 6, by contrast, shows CA increase from
    Session 1 --- before deep relationships have had time to form. A monotone increase
    from session 1 is more consistent with a mechanic that is active from session 1
    than with a relationship-formation process that takes multiple sessions to develop.

  \item \textbf{C4 (siege archetype, Supply Die mechanic):} Campaign 4 shows the
    monotone-rising trajectory under a different party archetype (siege-resource,
    not military-duty) and a different resource mechanic (Supply Die). The fact that
    a non-military archetype also produces a monotone-rising trajectory under a
    degrading-resource mechanic suggests the trajectory type is not military-archetype-specific.
\end{itemize}

Taken together, the evidence is consistent with the resource mechanic interpretation:
duty-adjacent parties without the mechanic (C5) show stable trajectories; the mechanic
appears in both a siege context (C4) and a military context (C6) and in both cases
produces the monotone-rising pattern.

\paragraph{Acknowledged limit.}
We cannot rule out that duty-bound party archetypes produce monotone CA increase
independent of resource mechanics. The C5 comparison is imperfect because the
professional-code archetype's emotional register is lower-affect by design (see
Section~4's discussion of C5's Route B designation). A closer comparison would be
a military-archetype campaign without any resource mechanic, which the corpus does not
contain. A controlled experiment varying resource mechanic presence/absence within
the same party archetype would resolve this.

\paragraph{Design recommendation.}
For practitioners, we recommend treating the $+3.0$ figure as an upper bound
attributable to the combination of unit-cohesion mechanics and duty-bound party
archetypes, and the $+0.6$ figure (C4, neutral siege party) as a lower bound
attributable to resource mechanics alone. The practitioner claim that is robust
across both estimates is that a legible degrading-resource mechanic accelerates
arrival at peak Character Agency by 3--4 sessions relative to no-resource baselines,
regardless of archetype. The magnitude of the advantage above that acceleration floor
may vary with party archetype.

\subsection{Limitations}

\paragraph{Archetype confound.}
As discussed in Section~\ref{subsec:archetype-confound}, the primary limitation is
the potential confound between resource mechanic and campaign archetype. A fully
controlled comparison would require two campaigns of the same archetype, one with and
one without a degrading-resource mechanic, with all other spine components held constant.
The Marathon corpus does not contain this comparison. The partial controls available
(C5, C7, C4) support the resource mechanic interpretation but cannot fully eliminate
the archetype alternative.

\paragraph{Single system.}
All 7 campaigns were run through the same system (Marathon D\&D Workshop, D\&D 5e
ruleset). The amplification effect may be specific to D\&D 5e's resource architecture
(spell slots, hit dice, action economy) and may not generalize to systems with
different base resource mechanics (e.g., Blades in the Dark's stress/trauma,
Ironsworn's supply/spirit tracks, PbtA moves). Replication in a different system
would be required to establish the finding as system-independent.

\paragraph{AI-simulated sessions.}
All sessions were AI-simulated playtests. Human player behavior may differ from
AI-simulated behavior in ways that affect resource management patterns. Specifically,
human players may show stronger or weaker preservation behavior than AI agents, and
human grief-paragraph engagement may not match the simulated pattern. The direction
of bias is unknown.

\paragraph{Small corpus.}
Two resource campaigns and five baselines is a small corpus for drawing generalizable
conclusions. The structural patterns (monotone-rising vs.\ late-rising) are consistent
across all 7 campaigns, but the quantitative advantage ($+0.6$, $+3.0$ panel-session
advantage) is derived from data that could be materially affected by a single unusual
campaign.
